\chapter{ADEQUAÇÃO ACADÊMICA E MANUTENÇÃO}
\label{chapter-F}
\section{Correspondência com o manual da FATEC}
\label{section-F-1}
A estrutura segue os elementos de TG II descritos no manual institucional de maio de 2024: capa, folha de rosto, resumo, abstract, listas pertinentes, sumário, introdução, fundamentação, método, resultados e discussão, conclusão e referências. Concepção e implementação foram desdobradas em capítulos próprios para documentar a engenharia do artefato \hyperref[ref:FATECMANUAL]{(FATEC CAMPINAS, 2024)}. O projeto não reproduz integralmente o layout gráfico do template.

A composição adota A4, fonte de corpo 12, títulos principais 14, margens superior/esquerda de 3 cm e inferior/direita de 2 cm, entrelinha 1,5 no corpo e recuo de 1,25 cm. Elementos como referências, legendas e resumos usam espaçamento simples. A contagem começa na folha de rosto, e os números aparecem a partir da introdução, no alto à direita. A matriz externa matriz-fatec.md permite conferir cada exigência e a localização.

Folha de aprovação, assinaturas e ficha catalográfica da versão final dependem dos responsáveis institucionais. O portal da FATEC orienta que a biblioteca fornece a ficha e que a aprovação é obtida com a orientação \hyperref[ref:FATECPORTAL]{(FATEC CAMPINAS, s.d.)}. Esses documentos não foram inventados nem preenchidos como aprovados. Dedicatória e agradecimentos são opcionais e devem refletir a manifestação pessoal da autora.

\section{Revisão e governança do documento}
\label{section-F-2}
A autora e a orientação devem revisar título, formulação dos objetivos, termos jurídicos, referências e declaração de assistência tecnológica. A defesa, eventual aprovação e alterações posteriores precisam de registros reais. O relatório é uma versão completa para revisão, não uma afirmação de submissão aceita. A conferência acadêmica não elimina a revisão jurídica indicada no aplicativo.

O conteúdo compartilhado em conteudo.md alimenta Word e LaTeX; os diagramas editáveis são gerados a partir de especificações locais. Atualizar primeiro a fonte e regenerar os artefatos evita versões divergentes. A versão LaTeX mantém capítulos, imagens e referências no mesmo projeto. O guia operacional de dez páginas acompanha o relatório, permitindo consulta rápida durante instalação e demonstração.
